\par\medskip
\noindent\textbf{Lema 6.2.}
Sea $A_0$ un $k$-álgebra local esencialmente de tipo finito tal que
\[
\widehat A_0\simeq k[[x_1,\ldots,x_n]]/(f),
\]
donde $f$ es una forma cuadrática no degenerada, y sea
$(R,\mathfrak X)$ una deformación versal formal de
$X_0=\operatorname{Spec}(A_0)$. Entonces:

\smallskip
\noindent\textbf{(1)} Los morfismos canónicos
\[
J_{\mathfrak X}(\mathfrak X|R)
 \longrightarrow J_R(\mathfrak X|R)
 \longrightarrow \operatorname{Spec}(\Lambda)
\]
son ambos isomorfismos.

\smallskip
\noindent\textbf{(2)} El ideal
$I_R(\mathfrak X|R)$ es principal y está generado por un generador
formal libre de $R$ sobre $\Lambda$; en otras palabras,
\[
I_R(\mathfrak X|R)=(t),
\qquad R=\Lambda[[t]].
\]

\smallskip
\noindent\textbf{(3)} El morfismo
$\mathfrak X\longrightarrow\operatorname{Spec}(\Lambda)$ es
formalmente liso.

\smallskip
\noindent\textit{Demostración.}
Escribamos $\mathfrak X=\operatorname{Spf}(A)$. Por 4.10(b),
$(R,\operatorname{Spf}(\widehat A))$ es una deformación versal formal
de $\widehat X_0=\operatorname{Spec}(\widehat A_0)$, donde
\[
\widehat A=\varprojlim_\nu A\otimes_RR_\nu.
\]
Ahora bien,
\[
\widehat A_0\simeq k[[x_1,\ldots,x_n]]/(f),
\qquad
\left(\frac{\partial f}{\partial x_1},\ldots,
      \frac{\partial f}{\partial x_n}\right)
=(x_1,\ldots,x_n).
\]
Se sigue inmediatamente que
\[
\bigl[\underline M_{\widehat X_0}(k[\varepsilon])\bigr]
\simeq k.
\]
Por consiguiente, por 4.14,
\[
R\simeq\Lambda[[t]],
\qquad
\widehat A\simeq R[[x_1,\ldots,x_n]]/(F-t),
\]
donde $F$ es el polinomio cuadrático de
$\Lambda[x_1,\ldots,x_n]\subset R[[x_1,\ldots,x_n]]$ obtenido de
$f$ levantando sus coeficientes a $\Lambda$.

La sucesión
\[
0\longrightarrow\widehat A
 \sisi{\overset{dF}{\relbar\joinrel\relbar\joinrel\relbar}}{\xrightarrow{\ dF\ }}
 \widehat A\,dx_1\oplus\cdots\oplus\widehat A\,dx_n
 \sisi{\relbar\joinrel\relbar\joinrel\relbar\joinrel\relbar\joinrel\relbar}{\longrightarrow}\widehat\Omega_{\widehat A|R}\longrightarrow0
\]
es exacta, donde
\[
1\longmapsto
 \frac{\partial F}{\partial x_1}dx_1+\cdots+
 \frac{\partial F}{\partial x_n}dx_n.
\]

%%%%%%%%%%  scan idx 137  |  folio 126  %%%%%%%%%%

Por tanto,
\[
I_{\widehat A}(\widehat A|R)
=\left(\frac{\partial F}{\partial x_1},\ldots,
       \frac{\partial F}{\partial x_n}\right)
=(x_1,\ldots,x_n),
\]
y por consiguiente
\[
\begin{aligned}
\widehat A/I_{\widehat A}(\widehat A|R)
&=\widehat A/(x_1,\ldots,x_n)\\
&=R[[x_1,\ldots,x_n]]/(x_1,\ldots,x_n,F-t)
 \simeq\Lambda.
\end{aligned}
\]
Así, el morfismo canónico
$\Lambda\longrightarrow
\widehat A/I_{\widehat A}(\widehat A|R)$ es un isomorfismo. La
comparación de las completaciones identifica este cociente con
$A/I_A(A|R)$ y, por consiguiente,
\[
\Lambda\xrightarrow{\ \sim\ }R/I_R(A|R).
\]
Esto demuestra (1) y (2). Finalmente, (3) se deduce de los
isomorfismos
\[
\begin{aligned}
\widehat A
&\simeq R[[x_1,\ldots,x_n]]/(F-t)\\
&\simeq\Lambda[[t,x_1,\ldots,x_n]]/(F-t)
 \simeq\Lambda[[x_1,\ldots,x_n]]
\end{aligned}
\]
de $\Lambda$-álgebras.
\hfill$\square$
